\documentclass[journal]{IEEEtran}
\usepackage{amsmath,amssymb,mathtools}
\usepackage{amsthm}
\usepackage{booktabs,array}
\usepackage{algorithm}
\usepackage{algpseudocode}
\usepackage{enumitem}
\usepackage{microtype}
\usepackage[hidelinks]{hyperref}
\newtheoremstyle{plainn}{}{}{\itshape}{}{\bfseries}{.}{.5em}%
  {\thmname{#1}\thmnumber{ #2}\thmnote{ (#3)}}
\theoremstyle{plainn}
\newtheorem{thm}{Theorem}
\newtheorem{prop}[thm]{Proposition}
\newtheorem{cor}[thm]{Corollary}
\newtheorem{claim}{Claim}
\theoremstyle{remark}
\newtheorem{remark}{Remark}
\renewcommand{\qedsymbol}{QED}

\begin{document}
\title{Minimum order and minimum-order construction for functional observers}
\author{}
\maketitle

\begin{abstract}
For a linear system with a prescribed block of functional rows, deciding whether some observer of a
given order exists is a question with a complete characterization, and its complexity is settled in a
companion paper. This paper asks the optimization question: what is the smallest order at which a
functional observer exists, and how is a minimum-order observer constructed? The feasible designs are
subspaces $W$ of a compressed row-pair space, constrained by $W_0\subseteq W\subseteq\mathcal T$ and
$\dim(W/W_0)\le r$, on which the two conditions are a row-space inclusion and the constant rank of a
pencil. Three results follow. First, the minimum order is $q_{\min}=r_0+t_{\min}$, where $t_{\min}$ is
the smallest quotient dimension of a feasible $W$, the structural bound $\omega+1\le\dim W\le2\omega$
gives $t_{\min}\le2\ell$ with $\ell=\operatorname{rank}(AN)$, and the feasible orders are monotone and
saturate, so the threshold problem is a single decision instance. Second, the spectral condition yields
the lower bound $q_{\min}\ge r_0+\max\bigl(0,\operatorname{rank}[U_0;V_0]-s_0+1\bigr)$, which rests on
a different necessary condition than the published algebraic bound $q_0$ and is strictly stronger on the
instance below, where $q_0=3$ and $q_{\min}=4$. Third, given a decision procedure, the minimum order is
found in $O(\log(\ell+2))$ calls and a minimum-order observer is recovered by lifting a basis of
$W/W_0$; the lifting is exact, while the algebraic solve is bounded in the arithmetic model, its bit
complexity being open.
\end{abstract}

\begin{IEEEkeywords}
functional observer; minimum order; subspace constraints; lower bounds; observer construction.
\end{IEEEkeywords}

\section{Introduction}\label{sec:intro}

Estimating a linear function of the state, rather than the whole state, is often what a control loop
needs, and the order of a functional observer can be far below that of a Luenberger observer for the
same system. The design question is quantitative: given the system and the functionals to be
estimated, what is the smallest number of observer rows for which a functional observer exists?

Two rank conditions characterize the observer rows, one algebraic and one spectral. The
characterization is complete in recent work: the design conditions are those of \cite{fer10,dar25}, and
\cite{fer26a,fer26b} determine, order by order, which orders are attainable, together with a
construction at the minimum order. Those results answer whether a given order is admissible and how to
design at it; they say less about the set of admissible orders as a set, and less about what happens
between the bound $q_0$ and the minimum order itself. Deciding the question for a prescribed order is
treated in a companion paper \cite{companion1} and is used here as an interface. Of the two conditions
from which a lower bound for the minimum order could be derived, only the algebraic one is used in that
literature: \eqref{eq:5a}
alone yields $q_0=\sum_j\eta_j$ from the functional observability indices, whereas the spectral
condition \eqref{eq:5b} is treated there as a test applied at each order, not as a source
of a bound. Two questions are therefore open. What is the structure of the set of feasible orders? What
separates the algebraic bound $q_0$ from the minimum order?

The answer is carried by a subspace: the free rows enter the two conditions only through the row space
$W$ of the compressed data, constrained by $W_0\subseteq W\subseteq\mathcal T$ and
$\dim(W/W_0)\le r$. A feasible design is thus a subspace with these constraints, the observer conditions
depend on $W$ alone, and the minimum-order question becomes a constrained dimension-minimization
problem; the structural bound $\omega+1\le\dim W\le2\omega$ of the companion paper is what makes the
search finite, by forcing the quotient dimension of a feasible $W$ to be at most $2\ell$. The paper
follows one chain---the structure and the standing results (Section~\ref{sec:standing}), monotonicity
and saturation and the formula $q_{\min}=r_0+t_{\min}$ (Section~\ref{sec:monotone}), the spectral lower
bound and its equality case (Section~\ref{sec:lowerbound}), and the interface that computes the order
and constructs the observer (Section~\ref{sec:compute})---and the conclusions are the following.

\begin{enumerate}
\item \emph{Minimum-order characterization} (Section~\ref{sec:monotone}). The feasible orders are
monotone and saturate at $r_0+\dim(\mathcal T/W_0)$, so the threshold problem is a single decision
instance, and
\begin{equation*}
\begin{gathered}
q_{\min}=r_0+t_{\min},\\
t_{\min}=\min\{\dim(W/W_0):W\text{ feasible}\}\le2\ell .
\end{gathered}
\end{equation*}
The three quantities play different roles: $r_0$ is fixed, $t_{\min}$ is minimized, and
$\ell=\operatorname{rank}(AN)$ bounds the search.
\item \emph{A spectral lower bound} (Section~\ref{sec:lowerbound}). With
$U_0=\operatorname{row}(L_0N)$, $V_0=\operatorname{row}(L_0AN)$ and $s_0=\dim W_0$,
\begin{equation*}
q_{\min}\ \ge\ r_0+\max\bigl(0,\ \operatorname{rank}[U_0;V_0]-s_0+1\bigr).
\end{equation*}
The published bound $q_0$ of \cite{fer26a,fer26b} comes from the algebraic condition; this one comes
from the spectral condition, so the two rest on different necessary conditions and the larger applies.
\item \emph{Computation and construction} (Section~\ref{sec:compute}). Given a decision procedure for
a prescribed order, the minimum order is found in $O(\log(\ell+2))$ calls, and a minimum-order observer is
recovered by lifting a basis of $W/W_0$ to free rows. The lifting is exact; the algebraic solve is
bounded in the arithmetic model, and its bit complexity is left open.
\item \emph{A published value corrected} (Appendix~\ref{app:example}). On the instance of
\cite[Sec.~VII]{fer26b} the reported minimum order is $5$; the lower bound above gives
$q_{\min}\ge4$ and an explicit four-row observer matrix satisfies both conditions, so
$q_{\min}=4$.
\end{enumerate}

Appendix~\ref{app:example} records the correction of the last item in full; the decision problem
itself, its finite certificate and its complexity bounds are not reproved here.

\section{Preliminaries}\label{sec:prelim}

\subsection{Notation and conventions}\label{sec:notation}

We take $C$ to have full row rank, $p=\operatorname{rank}C$: otherwise $C$ is replaced by a basis of
its row space, which leaves $\operatorname{row}[C;CA]$ unchanged and hence leaves \eqref{eq:5a} and
\eqref{eq:5b} unchanged, since both are rank identities on matrices whose rows are augmented by $C$ and
$CA$. The integer $p$ in the order range \eqref{eq:range} is this rank.

Consider the linear time-invariant system
\begin{equation*}
\dot{x}=Ax+Bu,\quad y=Cx,\quad z=L_0x,
\end{equation*}
where $A\in\mathbb Q^{n\times n}$, $C\in\mathbb Q^{p\times n}$ and
$L_0\in\mathbb Q^{r_0\times n}$.

Table~\ref{tab:notation} collects the notation. Let $q\ge r_0$ be the prescribed observer order,
with $q\le n$, and set $r=q-r_0$. The
observer rows are the rows of the block
\begin{equation*}
L=\begin{bmatrix}L_0\\R\end{bmatrix},\qquad R\in\mathbb R^{r\times n},
\end{equation*}
where $L_0$ is fixed while $R$ collects the $r$ free rows. Throughout we call the $r_0$ rows of
$L_0$ the \emph{prescribed rows}, the $r$ rows of $R$ the \emph{free rows}, and the $q$ rows of $L$
the \emph{observer rows}; only real observers are considered.

All given rational data are encoded in binary, and $\mathcal L$ denotes the total input
bit length. The free parameter $R$ is not part of the given input data and is therefore
not counted in $\mathcal L$. State vectors are column vectors, and matrices act on them by
the usual left multiplication.

Let $P(\lambda)$ be a rational matrix pencil of generic rank $\omega$, and call its
$\omega\times\omega$ minors maximal minors. Ranks at spectral points $\lambda\in\mathbb C$
are always computed over $\mathbb C$. Because the maximal minors have rational
coefficients, their gcd is a nonzero constant in $\mathbb C[\lambda]$ exactly when it is a
nonzero constant in $\mathbb Q[\lambda]$; this equivalence is used without further
mention.

\subsection{The observer conditions}\label{sec:conditions}

A functional observer of order $q$ is a matrix $L=\begin{bmatrix}L_0 \\ R\end{bmatrix}\in\mathbb R^{q\times n}$ such
that
\begin{equation}\label{eq:5a}
\operatorname{rank}\begin{bmatrix}LA\\ L\\ CA\\ C\end{bmatrix}
=\operatorname{rank}\begin{bmatrix}L\\ CA\\ C\end{bmatrix},
\end{equation}
\begin{equation}\label{eq:5b}
\operatorname{rank}\begin{bmatrix}\lambda L-LA\\ CA\\ C\end{bmatrix}
=\operatorname{rank}\begin{bmatrix}L\\ CA\\ C\end{bmatrix},
\quad \forall\lambda\in\mathbb C.
\end{equation}
These are the Fernando--Trinh--Jennings conditions, numbered (5a) and (5b) in
\cite{fer10}. Condition~\eqref{eq:5a} is algebraic; \eqref{eq:5b} is spectral, and it is
required at every $\lambda\in\mathbb C$, in particular at all finite eigenvalues of $A$.

\subsection{Decision, optimization and construction}\label{sec:problems}

Three problems are posed for \eqref{eq:5a} and \eqref{eq:5b}, in increasing order of
demand.

\emph{Decision.} $\mathrm{DEC}$ takes $(A,C,L_0)$ and $q$ and asks whether some $R$
satisfies \eqref{eq:5a} and \eqref{eq:5b}; no witness is required. This is the problem
studied in this paper.

\emph{Optimization.} $\mathrm{OPT}$ asks for the minimum feasible order
\begin{equation*}
q_{\min}=\min\{q:\mathrm{DEC}(q)\},
\end{equation*}
where $\mathrm{DEC}(q)$ abbreviates $\mathrm{DEC}(A,C,L_0,q)$.

\emph{Construction.} $\mathrm{FIND}$ asks for a witness $R$ at the minimum feasible
order, $q=q_{\min}$.

The threshold problem is
\begin{equation}\label{eq:mindec}
\mathrm{MINDEC}(Q)=\bigvee_{q=r_0}^{\min(Q,n)}\mathrm{DEC}(q),
\end{equation}
a finite disjunction of decision instances. This paper studies $\mathrm{OPT}$ and
$\mathrm{FIND}$; $\mathrm{DEC}$ itself, its finite certificate and its complexity bounds are treated in
the companion paper, whose decision procedure we use as an interface. What is added here is the
structure of the feasible orders: they are monotone in $q$ and saturate at $r_0+m$, which reduces
\eqref{eq:mindec} to a single instance and confines $\mathrm{OPT}$ to at most $2\ell+1$ values of
$t=q-r_0$.

\begin{table}[!t]
\centering
\caption{Notation. The first block is fixed by the input, the second by the compression, the third by
the subspace reduction of Section~\ref{sec:standing}.}
\label{tab:notation}
\footnotesize
\begin{tabular}{@{}ll@{}}
\toprule
Symbol & Meaning \\
\midrule
$K_0=\ker[C;CA]$ & common kernel of $C$ and $CA$ \\
$k=\dim K_0$ & its dimension \\
$N$ & basis matrix of $K_0$ \\
$M=[N\ AN]$ & compressed row-pair matrix \\
\midrule
$H=LN$, $H_A=LAN$ & compressed observer blocks \\
$\omega=\operatorname{rank}H$ & compressed order \\
$P(\lambda)=\lambda H-H_A$ & compressed pencil \\
\midrule
$\mathcal T=\operatorname{row}M$ & attainable row-pair space \\
$d_0=\dim\mathcal T$ & its dimension \\
$\kappa_{\rm obs}=r\,d_0$ & number of free-row coordinates \\
$S=\operatorname{row}(AN)\subseteq\mathbb R^k$ & space of the second blocks \\
$\ell=\operatorname{rank}(AN)\le k$ & its dimension; the FPT parameter \\
$W=\operatorname{row}[H\ H_A]$ & row space of the compressed pairs \\
$W_0=\operatorname{row}(L_0M)$ & contribution of the prescribed rows \\
$s_0=\dim W_0$, $m=d_0-s_0$ & dimensions of $W_0$ and of $\mathcal T/W_0$ \\
$\mathcal V=\mathcal T\cap(S\times S)$ & admissible ambient space \\
$m'=\dim\mathcal V-s_0$ & chart dimension \\
$\pi_1,\pi_2,\pi_\lambda$ & projections of $\mathbb R^{2k}$ \\
\bottomrule
\end{tabular}
\end{table}

Two published bounds confine the orders that have to be examined. Let $q_0$ be the algebraic
lower bound of \cite{fer26a,fer26b}: the smallest order at which \eqref{eq:5a} can be satisfied,
computed from the functional observability indices by a greedy rank recursion, hence in polynomial
time. Every feasible order satisfies
\begin{equation}\label{eq:range}
q_0\le q\le n_0-p,\qquad n_0=\operatorname{rank}\bigl[C;CA;\dots;CA^{n-1}\bigr],
\end{equation}
see \cite{fer26a,fer26b}; for an observable pair $n_0=n$. The upper end is not merely a bound in the
row-slot model: by Lemma~\ref{lem:monotone} the predicate is monotone in $q$, so every order at or
above $n_0-p$ is feasible as soon as that one is, and no information is lost by stopping there. In
terms of the quotient dimension the range \eqref{eq:range} reads
$q_0-r_0\le t\le n_0-p-r_0$, and the reduction of Section~\ref{sec:standing} adds the independent
bound $t\le2\ell$ from Proposition~\ref{prop:minorder}.

\subsection{Standing results from the companion paper}\label{sec:standing}

This subsection restates, without proof, the setting and the four results of the companion paper
\cite{companion1} on which everything below rests.

\begin{claim}[Compression and pencil form]\label{claim:compress}
Let
\begin{equation*}
K_0=\ker\begin{bmatrix}C\\ CA\end{bmatrix},\qquad k=\dim K_0,
\end{equation*}
and let $N\in\mathbb R^{n\times k}$ be a basis matrix for $K_0$. Define
\begin{equation*}
H=LN\in\mathbb R^{q\times k},\quad H_A=LAN\in\mathbb R^{q\times k},\quad \omega=\operatorname{rank}H .
\end{equation*}
Then \eqref{eq:5a} holds if and only if $\operatorname{row}H_A\subseteq\operatorname{row}H$. If
\eqref{eq:5a} holds and $F$ is such that $H_A=FH$, then \eqref{eq:5b} holds if and only if
\begin{equation}\label{eq:pencil}
\operatorname{rank}(\lambda H-H_A)=\omega\qquad\text{for every }\lambda\in\mathbb C,
\end{equation}
and $(\lambda H-H_A)=(\lambda I-F)H$ shows that $\operatorname{rank}P(\lambda)\le\omega$ always, so
\eqref{eq:pencil} asks for equality; equivalently, the $\omega\times\omega$ maximal minors of
$P(\lambda)=\lambda H-H_A$ have no common complex root \cite[Ch.~IV]{lan02}. If \eqref{eq:5a} and
\eqref{eq:5b} both hold, then $\operatorname{rank}H_A=\operatorname{rank}H$: \eqref{eq:5a} gives
$\operatorname{row}H_A\subseteq\operatorname{row}H$ and \eqref{eq:pencil} at $\lambda=0$ gives
$\operatorname{rank}H_A=\operatorname{rank}H$. (Compression and pencil criterion, Claim~2 of
\cite{companion1}, which rests on the rank identity for a stacked matrix \cite[Ch.~0]{hj13}.)
\end{claim}

\begin{claim}[Attainable row pairs]\label{claim:rows}
Each free row $r_i\in\mathbb R^{1\times n}$ of $R$ determines the row pair
\begin{equation*}
(H_i,(H_A)_i)=r_i\,[N\ \ AN] \in\mathbb R^{1\times 2k}.
\end{equation*}
This attainable pair is uniquely specified by its coordinates relative to a basis of the subspace of
attainable pairs
\begin{equation*}
\begin{gathered}
\mathcal T=\operatorname{row}[N\ \ AN]\subseteq\mathbb R^{1\times 2k},\\
d_0=\dim\mathcal T=\operatorname{rank}[N\ \ AN];
\end{gathered}
\end{equation*}
a rational basis of $\mathcal T$ therefore allows each attainable pair generated by a free row to be
parameterized by exactly $d_0$ real coordinates, which describes the $r$ free rows by
\begin{equation*}
\kappa_{\rm obs}=r\,d_0
\end{equation*}
real coordinates, and $k\le d_0$ because $N$ has full column rank. (Attainable row pairs and the
effective dimension, Claim~3 of \cite{companion1}.)
\end{claim}

The count $\kappa_{\rm obs}$ is not the parameter of the algorithms: it grows with $r$, whereas the
parameter that enters Theorem~\ref{thm:interface} is $\ell\le k$.

\begin{remark}[Basis invariance]
If $N\mapsto NQ$ with nonsingular $Q\in\mathbb R^{k\times k}$, then
$[N\ \ AN]\mapsto[N\ \ AN]\operatorname{diag}(Q,Q)$ with $\operatorname{diag}(Q,Q)$ nonsingular, so
$\mathcal T$, $d_0$ and $\kappa_{\rm obs}=rd_0$ are unchanged; the free-row coordinates undergo the
nonsingular change $r_i\mapsto r_i\operatorname{diag}(Q^{-1},Q^{-1})$.
\end{remark}

\begin{claim}[Sandwich and the subspace reading]\label{claim:sandwich}\label{claim:W}
The free rows enter the two conditions only through the row space of the compressed pair. Writing the
row pair of a free row as a row of $M=[N\ AN]$, let
\begin{equation*}
W=\operatorname{row}[H\ H_A]\subseteq\mathbb R^{2k},
\end{equation*}
so that $W_0=\operatorname{row}(L_0M)\subseteq W\subseteq\mathcal T=\operatorname{row}M$. Let
\begin{equation}\label{eq:ell}
\begin{gathered}
d=\dim\mathcal T,\qquad s_0=\dim W_0,\qquad m=d-s_0,\\
\ell=\operatorname{rank}(AN)\le k,\qquad S=\operatorname{row}(AN)\subseteq\mathbb R^k ,
\end{gathered}
\end{equation}
and let $\pi_1(h,g)=h$, $\pi_2(h,g)=g$, $\pi_\lambda(h,g)=\lambda h-g$ be the projections of
$\mathbb R^{2k}=\mathbb R^k\times\mathbb R^k$. Because \eqref{eq:5b} quantifies over complex
$\lambda$, the projection $\pi_\lambda$ is applied to the complexification
$W_{\mathbb C}=W\otimes_{\mathbb R}\mathbb C\subseteq\mathbb C^{2k}$, and dimensions are over
$\mathbb C$; a real matrix has the same rank over $\mathbb R$ and over $\mathbb C$. Then:

\emph{(i)} for $L=[L_0;R]$ one has $W_0\subseteq W\subseteq\mathcal T$ and $\dim(W/W_0)\le r$, and
conversely every subspace $W$ with these three properties is realized by a free block $R$ with exactly
$t=\dim(W/W_0)$ rows, so $\mathrm{DEC}(q)$ holds if and only if some such $W$ satisfies the two
conditions in (ii);

\emph{(ii)} condition \eqref{eq:5a} holds if and only if $\pi_2(W)\subseteq\pi_1(W)$, and condition
\eqref{eq:5b} holds if and only if $\dim\pi_\lambda(W)=\dim\pi_1(W)$ for every $\lambda\in\mathbb C$.
(Sandwich and realizability, Claim~4, and the criterion on $W$, Claim~5 of \cite{companion1}.)
\end{claim}

\begin{claim}[Structural bound]\label{claim:struct}
Let $W$ be feasible and $\omega=\dim\pi_1(W)>0$. Then
\begin{equation}\label{eq:struct}
\omega+1\le\dim W\le2\omega .
\end{equation}
A feasible $W$ therefore satisfies $\pi_1(W)=\pi_2(W)=\operatorname{row}H_A$, and since
$\operatorname{row}(LAN)\subseteq\operatorname{row}(AN)=S$, its projections lie in $S$. (Structural
bound, Claim~6 of \cite{companion1}; the projection equality is Claim~\ref{claim:compress} above.)
\end{claim}

The rest of the paper uses \eqref{eq:pencil} rather than \eqref{eq:5b}.



\section{The minimum order and its spectral bound}\label{sec:minorder}

Two things are proved here: that the set of feasible orders has a simple shape, and that the spectral
condition alone yields a lower bound on its minimum.

\subsection{Monotonicity, saturation, and the minimum order}\label{sec:monotone}

\begin{claim}[Monotonicity and saturation]\label{lem:monotone}
$\mathrm{DEC}(q)\Rightarrow\mathrm{DEC}(q+1)$, and $\mathrm{DEC}(r_0+r)=\mathrm{DEC}(r_0+m)$
whenever $r\ge m$. Consequently $\mathrm{MINDEC}(Q)=\mathrm{DEC}(\min(Q,n))$ for
$Q\ge r_0$.
\end{claim}

\begin{proof}
Adjoining a zero row to $R$ leaves $H$, $H_A$, and the ranks in \eqref{eq:5a} and
\eqref{eq:5b} unchanged, which gives monotonicity. Saturation follows from
$\dim(W/W_0)\le\dim(\mathcal T/W_0)=m$ and Claim~\ref{claim:sandwich}: any feasible $W$ with
$\dim(W/W_0)\le r$ is already realized with at most $m$ free rows. The disjunction in
\eqref{eq:mindec} then reduces to its largest term.
\end{proof}

\begin{prop}[Minimum order]\label{prop:minorder}
Let $t_{\min}=\min\{\dim(W/W_0):W \text{ feasible}\}$, and $t_{\min}=q_{\min}=+\infty$ when no $W$
is feasible. Then $q_{\min}=r_0+t_{\min}$ and $t_{\min}\le2\ell$.
\end{prop}

\begin{proof}
Claim~\ref{claim:sandwich} identifies feasibility with the existence of a feasible $W$ of
quotient dimension $\le r$, and Lemma~\ref{lem:monotone} shows that the smallest such $r$ is
$t_{\min}$. For the bound, a feasible $W$ satisfies $W\subseteq U\times U$ with $U=\pi_1(W)=\pi_2(W)$, and by
Claim~\ref{claim:compress} $U=\operatorname{row}H_A\subseteq S$; hence $\dim U=\omega\le\ell$ and
$\dim W\le2\omega\le2\ell$, that is $t_{\min}\le\dim W-s_0\le2\ell$.
\end{proof}

\subsection{The spectral lower bound, and when it is attained}\label{sec:lowerbound}

\begin{prop}[A computable lower bound]\label{prop:lowerbound}
Let $U_0=\operatorname{row}(L_0N)$ and $V_0=\operatorname{row}(L_0AN)$, and let $W$ be feasible with
$t=\dim(W/W_0)$ and $\omega=\dim\pi_1(W)$. Then
\begin{equation}\label{eq:lb}
\operatorname{rank}\,[\,U_0;\,V_0\,]\ \le\ \omega\ \le\ s_0+t-1,
\end{equation}
hence $t\ge\operatorname{rank}[U_0;V_0]-s_0+1$ and
\begin{equation}\label{eq:qmin-lb}
q_{\min}\ \ge\ r_0+\max\bigl(0,\ \operatorname{rank}[U_0;V_0]-s_0+1\bigr).
\end{equation}
\end{prop}

\begin{proof}
Since $W\supseteq W_0$ the projections satisfy $\pi_1(W)\supseteq\pi_1(W_0)=U_0$ and
$\pi_2(W)\supseteq\pi_2(W_0)=V_0$. By Claim~\ref{claim:W} at $\lambda=0$ we have
$\pi_1(W)=\pi_2(W)=U$ with $\dim U=\omega$, so $U\supseteq U_0+V_0$ and
$\omega\ge\operatorname{rank}[U_0;V_0]$. The upper bound is Claim~\ref{claim:struct}:
$\omega+1\le\dim W=s_0+t$. Combining and using $q_{\min}=r_0+t_{\min}$ with $t_{\min}\ge0$ gives
\eqref{eq:qmin-lb}.
\end{proof}

Proposition~\ref{prop:lowerbound} is the spectral counterpart of the algebraic bound of
\cite{fer26a,fer26b}. That bound, $q_0=\sum_j\eta_j$, limits the orders at which \eqref{eq:5a} can
hold; \eqref{eq:lb} limits the orders at which \eqref{eq:5b} can hold, since it uses only the condition
at $\lambda=0$ and the structural bound. The two are independent, both are computed in polynomial time,
and the larger one applies; on the instance of Appendix~\ref{app:example} the bound gives $4$ while
$q_0=3$. When $\operatorname{rank}[U_0;V_0]\le s_0$ the bound is vacuous while $q_0$ may be positive, so
neither bound dominates the other. To our knowledge the literature states no lower bound derived from the spectral condition: the
design conditions of \cite{fer10,dar25,fer26a,fer26b} and the order bound $q_0$ all concern
\eqref{eq:5a}, and the spectral condition is treated there as a test applied at each order rather than
as a source of a bound.

\begin{prop}[When the bound is attained]\label{prop:tight}
Let $U_0=\pi_1(W_0)$, $V_0=\pi_2(W_0)$, $a=\operatorname{rank}[U_0;V_0]$ and $a-s_0+1>0$. Then
$t_{\min}=a-s_0+1$ if and only if there exist $b\in U_0+V_0$, $b\neq0$, and a linear map
$G:U_0+V_0\to U_0+V_0$ such that
\begin{equation*}
W_0\subseteq W(G,b):=\bigl\{(h,Gh+\alpha b):h\in U_0+V_0,\ \alpha\in\mathbb R\bigr\}\subseteq\mathcal T
\end{equation*}
and $\operatorname{rank}\bigl[\lambda I-G\ \ b\bigr]=a$ for every $\lambda\in\mathbb C$.
\end{prop}

\begin{proof}
If such a pair $(G,b)$ exists, then $W(G,b)$ is a subspace of $\mathcal T$ with $W_0\subseteq W(G,b)$,
$\dim W(G,b)=a+1$ because $b\neq0$, and $\pi_\lambda(W(G,b))=(\lambda I-G)(U_0+V_0)+\operatorname{span}\{b\}$
has dimension $a$ for every $\lambda$; by Claim~\ref{claim:W} it is feasible, so its quotient
dimension $a+1-s_0$ bounds $t_{\min}$ above and \eqref{eq:qmin-lb} bounds it below.

Conversely, let $W$ attain $t_{\min}=a-s_0+1$, so $\dim W=s_0+t_{\min}=a+1$. By
Claim~\ref{claim:W} $\pi_1(W)=\pi_2(W)=:U$; since $\pi_1(W)\supseteq U_0$ and
$\pi_2(W)\supseteq V_0$ we have $U\supseteq U_0+V_0$, and $\dim U\ge a$ while
$\dim U\le\dim W-1=a$ by Claim~\ref{claim:struct}; hence $U=U_0+V_0$ and $\dim U=a$. The kernel of
$\pi_1|_W$ is then one-dimensional, say $\operatorname{span}\{(0,b)\}$ with $b\neq0$, and choosing a
section of $\pi_1|_W$ gives a linear $G:U\to U$ with $W=W(G,b)$. Feasibility gives
$W_0\subseteq W\subseteq\mathcal T$, and the spectral condition gives
$\dim\bigl((\lambda I-G)U+\operatorname{span}\{b\}\bigr)=a$, which is
$\operatorname{rank}[\lambda I-G\ \ b]=a$.
\end{proof}

\begin{remark}[On tightness]\label{rem:lbtight}
Proposition~\ref{prop:tight} is the exact criterion: the bound is attained exactly when the
prescribed subspace admits a one-dimensional extension $W(G,b)$ inside $\mathcal T$ whose extra
direction makes the single-input pair $(G,b)$ controllable. This is a genuine additional requirement,
and it does not yield a polynomial-time tightness test. The bound \eqref{eq:qmin-lb} is attained on the
instance of Appendix~\ref{app:example}, where it gives
$4$ and $q_{\min}=4$. We do not claim that it is tight in general: like the algebraic bound $q_0$ it is
a lower bound, and the difference
$q_{\min}-\bigl(r_0+\max(0,\operatorname{rank}[U_0;V_0]-s_0+1)\bigr)$ can be positive. What it does
provide is a test for the orders below it, which then need not be examined at all.
\end{remark}

Both statements are turned into procedures in Section~\ref{sec:compute}: monotonicity makes the
threshold problem a single decision instance, and \eqref{eq:qmin-lb} removes the orders below it from
the search.

\section{Computing and constructing}\label{sec:compute}

Both problems are reduced to the decision problem through an explicit interface, so that no complexity
argument of the companion paper has to be repeated here. The interface is stated once, in the section
body, and both subsections rest on it: the first searches the order, the second lifts a witness at the
order found.

\begin{thm}[Interface]\label{thm:interface}
Suppose that $\mathrm{DEC}(q)$ can be decided in time $T(q)$ for every order $q$ in the window
\eqref{eq:range}, and let $\ell=\operatorname{rank}(AN)$. Then:
\emph{(i)} $t_{\min}$, and hence $q_{\min}$, is computed with $O(\log(\ell+2))$ calls to $\mathrm{DEC}$;
\emph{(ii)} given an accepted order and a subspace $W$ realizing it, an observer matrix $R$ with exactly
$t=\dim(W/W_0)$ rows and $R[N\ AN]=Z$ is obtained by solving $t$ rational linear systems of size
$n\times2k$.
\end{thm}

\begin{proof}
(i) By Proposition~\ref{prop:minorder} the feasible orders are monotone in $q$ and $t_{\min}\le2\ell$,
so a binary search over the window \eqref{eq:range} locates the least feasible order in $O(\log\ell)$
calls. (ii) By Claim~\ref{claim:sandwich} a basis $w_1,\dots,w_t$ of $W/W_0$ lifts to rows $r_i$ with
$r_i[N\ AN]=w_i$, and padding with zero rows realizes $W$ with exactly $t$ rows.
\end{proof}

Given the decision procedure of the companion paper, whose cost is
$f(\ell)\operatorname{poly}(\mathcal L)$ bit operations, the interface gives $\mathrm{OPT}$ in
$f(\ell)\operatorname{poly}(\mathcal L)$ bit operations and a construction bounded in the arithmetic
model; the statements of the two procedures follow.

\subsection{Minimum order by binary search}\label{sec:bsearch}

\begin{cor}[$\mathrm{OPT}$]\label{cor:opt}
$\mathrm{OPT}$ is decidable in $\exp\bigl(O(\ell^2\log(\ell+2))\bigr)\operatorname{poly}(\mathcal L)$
bit operations.
\end{cor}

\begin{proof}
By Proposition~\ref{prop:minorder} $q_{\min}=r_0+t_{\min}$, and by \eqref{eq:range} and the same
proposition $t_{\min}$ lies in $[\max(0,q_0-r_0),\,\min(Q-r_0,\,n_0-p-r_0,\,2\ell)]$; that is the
initial interval of Algorithm~\ref{alg:opt}, so the search is confined before it starts. By
Lemma~\ref{lem:monotone} the predicate ``$\mathrm{DEC}(r_0+t)$ holds'' is monotone in $t$, so the
algorithm maintains the invariant that the least feasible $t$ lies in $[\mathit{lo},\mathit{hi}]$,
halves the interval at each step, and returns $t_{\min}$ after $O(\log(\ell+2))$ calls to the decision
procedure of Theorem~\ref{thm:interface}, the logarithmic factor being absorbed by the cost of that
procedure. The
test at the head is what distinguishes ``no feasible order up to $Q$'' from a finite answer, and it
is correct because the predicate is upward closed.
\end{proof}
Algorithm~\ref{alg:opt} makes the search explicit: $\mathrm{DEC}$ is used as a subroutine, and the
monotonicity of Lemma~\ref{lem:monotone} is what makes a binary search correct.


\begin{algorithm}[!t]
\caption{Minimum order $\mathrm{OPT}$, with $\mathrm{DEC}$ as a subroutine}
\label{alg:opt}
\begin{algorithmic}[1]
\Require $A,C,L_0,Q$; \ $\ell=\operatorname{rank}(AN)$; \ $r_0=\operatorname{rows}(L_0)$;
$q_0$ as in \eqref{eq:range}.
\Ensure $q_{\min}$, or $+\infty$ if no order up to $Q$ is feasible.
\State $\mathit{lo}\gets\max(0,\,q_0-r_0)$
\State $\mathit{hi}\gets\min\bigl(Q-r_0,\ n_0-p-r_0,\ 2\ell\bigr)$
\If{$\mathit{lo}>\mathit{hi}$ \textbf{or} $\mathrm{DEC}(r_0+\mathit{hi})$ returns \textsc{infeasible}}
  \State \Return $+\infty$
\EndIf
\While{$\mathit{lo}<\mathit{hi}$}
  \State $\mathit{mid}\gets\lfloor(\mathit{lo}+\mathit{hi})/2\rfloor$
  \If{$\mathrm{DEC}(r_0+\mathit{mid})$ returns \textsc{feasible}}
    \State $\mathit{hi}\gets\mathit{mid}$
  \Else
    \State $\mathit{lo}\gets\mathit{mid}+1$
  \EndIf
\EndWhile
\State \Return $r_0+\mathit{lo}$
\end{algorithmic}
\end{algorithm}



\begin{remark}[Room for improvement]\label{rem:refine}
The complexity bound does not depend on how tightly the window is chosen, so
Algorithm~\ref{alg:opt} uses only the published lower bound $q_0$ and the structural bound $2\ell$.
Two cheap refinements are available and are left out on purpose, to keep the algorithm free of
instance-specific preprocessing. First, \eqref{eq:qmin-lb} gives a lower bound that is computed in
polynomial time from $U_0$ and $V_0$ and can be strictly larger than $q_0$; on the instance of
Appendix~\ref{app:example} it equals $4$ while $q_0=3$. Replacing $\max(0,q_0-r_0)$ by the larger of
the two bounds in line~2 shrinks the window without affecting correctness. Second, the structural
bound \eqref{eq:struct} also confines $\omega$ on each chart, which prunes branches of the inner loop
in the algorithm of the companion paper. Neither refinement changes any asymptotic statement, and
neither is
needed for the results above.
\end{remark}

\subsection{Constructing a minimum-order observer}\label{sec:construct}

Procedure~\ref{proc:find} states the construction; step~1 is the only step that is not elementary,
and it is a citation rather than a description.

\paragraph*{Procedure FIND (arithmetic model)}\label{proc:find}
\begin{enumerate}
\item Solve the system $\Psi_{t,\omega}$ of the companion paper by the sampling-point algorithm of
\cite[Ch.~13]{bpr06}, for the least $t$ accepted by $\mathrm{DEC}$ and for some $\omega$ that
succeeds. This returns $Y$ and $x$ as exact real algebraic numbers, with the arithmetic bound of
\cite[Ch.~13]{bpr06}. We neither describe that algorithm nor bound its bit cost
(Remark~\ref{rem:find-bit}).
\item Form $W=\operatorname{row}[XB\ \ GB]$ from the solution of $\Psi_{t,\omega}$, choose a basis
$w_1,\dots,w_t$ of $W/W_0$, and solve the
linear systems $u_iM=w_i$ for $i=1,\dots,t$, whose coefficient matrix $M$ is rational while the
right-hand side is algebraic; this costs $\operatorname{poly}(\mathcal L)$ arithmetic operations in the
arithmetic model, as the bit-size estimates of the companion paper show.
\item Return $R=[u_1;\dots;u_t;0]$ padded with $Q-r_0-t$ zero rows.
\end{enumerate}

\paragraph*{The witness system $\Psi_{t,\omega}$}\label{sec:psi}
Step~1 of the procedure takes as given a subspace $W$ realizing an accepted order. This paragraph
records the algebraic system that produces such a $W$, so that the construction can be read without the
companion paper. Nothing is re-proved: the system is stated in the notation of this paper, and the
companion paper contains the derivation and the bound on its size.

Fix an order $q=r_0+t$ and a rank $\omega\le\min(q,k)$, and let $J\subseteq[k]$ with $|J|=\omega$.
The unknown chart matrices are
\begin{equation*}
X,G\in\mathbb R^{q\times\omega},\qquad x\in\mathbb R^{2\omega},
\end{equation*}
and $B\in\mathbb R^{\omega\times k}$ is the matrix determined by $X$ and $J$ through $XB=X$ on the
columns in $J$ and $B_{:,J}=I_\omega$; this is possible exactly when the columns of $X_J$ are
independent, which the system enforces below. Then $W=\operatorname{row}[XB\ \ GB]$, and
$\Psi_{t,\omega}$ is the disjunction of the following systems over $J$, with $h^0=L_0N$,
$g^0=L_0AN$ and $D$ a rational matrix with $\ker D=\operatorname{row}[N\ AN]$:

\begin{enumerate}\itemsep2pt
\item $e_\omega(X^{\mathsf T}X)>0$, the elementary symmetric polynomial of degree $\omega$ in the
eigenvalues of $X^{\mathsf T}X$ being positive: this is $\operatorname{rank}X=\omega$, hence
$\dim W=\omega$.
\item the rank-consistency equations $\mathcal M(X,G)\,b=c(X,G)$ of the companion paper, which say
that $B$ is defined by $X$ alone and is consistent with the second block $G$;
\item $X_iB=h_i^0$ and $G_iB=g_i^0$ for $i\le r_0$, the prescribed rows;
\item $D\,[(X_iB)^{\mathsf T};(G_iB)^{\mathsf T}]=0$ for $i>r_0$, which places the remaining row pairs
in $\mathcal T$;
\item $Z^{\mathsf T}Zx=e_0$, where $Z$ is the matrix of coefficient windows of the $\omega\times\omega$
maximal minors of $\lambda X-G$ and $e_0$ is the first elementary symmetric polynomial: this is the
finite form of \eqref{eq:pencil}.
\end{enumerate}

Items (i)--(ii) make $W$ a well-defined $\omega$-dimensional subspace of the chart, (iii)--(iv) place it
in the sandwich $W_0\subseteq W\subseteq\mathcal T$ with $\dim(W/W_0)=t$, and (v) is the spectral
condition by the companion paper's certificate. The system is thus the compressed form of
\eqref{eq:5a}--\eqref{eq:5b} at order $q=r_0+t$, and its real solutions are exactly the feasible $W$
of that order with $\dim\pi_1(W)=\omega$.

Two points are worth isolating, because the construction above uses them.
First, a real solution exists if and only if $\mathrm{DEC}(r_0+t)$ holds with $\omega$ as the rank of
the compressed block; that is the equivalence proved in the companion paper. Second, the solution set
is a real algebraic set defined by $O(\ell^2)$ polynomials in $\ell^2+2\ell$ variables of degree
$O(\ell)$, so a real sample point exists whose coordinates are algebraic numbers of degree
$D^{O(N)}=\exp\bigl(O(\ell^2\log(\ell+2))\bigr)$ --- the bound quoted in
Corollary~\ref{cor:find}. Reducing the coordinates to the row space $W$ is the step performed in the
arithmetic model.

\begin{cor}[$\mathrm{FIND}$, arithmetic model]\label{cor:find}
If the instance is feasible, a witness $R$ can be computed in
$\exp\bigl(O(\ell^2\log(\ell+2))\bigr)\operatorname{poly}(\mathcal L)$ arithmetic operations, as
exact real algebraic numbers whose number and degree are $f(\ell)$.
\end{cor}

\begin{proof}[Proof of Corollary~\ref{cor:find}]
Step~1 is the sampling-point bound of \cite[Ch.~13]{bpr06}: a nonempty real algebraic set defined by
$s$ polynomials in $N$ variables of degree at most $D$ and coefficient bit length $\beta$ has a sample
point whose coordinates are algebraic numbers of degree $D^{O(N)}$ with isolating intervals of bit
length $\beta D^{O(N)}$; with the parameter counts of the companion paper this is
$f(\ell)\operatorname{poly}(\mathcal L)$ arithmetic
operations and $f(\ell)\operatorname{poly}(\mathcal L)$ output bits. Steps~2 and~3 are
Claim~\ref{claim:sandwich}: a basis of $W/W_0$ lifts to free rows by solving $uM=w$, and padding with
zero rows realizes $W$ with exactly $t$ free rows. Since $W$ satisfies the two conditions
(Claim~\ref{claim:W}), the resulting $L=[L_0;R]$ does too, and its order is $r_0+t$.
\end{proof}

\begin{remark}[What is not claimed]\label{rem:find-bit}
Corollary~\ref{cor:find} is an arithmetic bound on the construction together with a bound on the
output length. A bound in the bit model for producing the witness is not asserted: the ordinary
analysis of \cite{bpr06} gives arithmetic operations only, and their Chapter~8 states explicitly
that bit-length estimates do not follow from the arithmetic analysis. We therefore state
$\mathrm{FIND}$ in the arithmetic model and leave its bit complexity open
(Section~\ref{sec:limits}).
\end{remark}

Deciding the minimum order is bounded in the bit model; the construction is bounded only in the
arithmetic model, and its bit complexity is open. Both statements are conditional on the predicate
$\mathrm{DEC}$ of Section~\ref{sec:problems}: the bound of Corollary~\ref{cor:opt} is the bound of the
companion paper's decision procedure multiplied by the $O(\log(\ell+2))$ calls of
Theorem~\ref{thm:interface}, and the arithmetic bound of Corollary~\ref{cor:find} inherits the
sampling-point cost of \cite{bpr06}. No bound on $\mathrm{OPT}$ is claimed independently of that
predicate.

\section{Discussion and conclusion}\label{sec:conclusions}

The feasible designs form a constrained family of subspaces, $W_0\subseteq W\subseteq\mathcal T$ with
$\dim(W/W_0)\le r$, and everything in this paper follows from that reading. The family is monotone and
saturates, so the threshold problem is a single decision instance and the minimum order is
$q_{\min}=r_0+t_{\min}$ with $t_{\min}\le2\ell$ (Proposition~\ref{prop:minorder}). The spectral
condition gives the lower bound \eqref{eq:qmin-lb}, which rests on a different necessary condition than
the algebraic bound $q_0$ of \cite{fer26a,fer26b} and is strictly stronger on the instance of
Appendix~\ref{app:example}. Given a decision procedure for a prescribed order, minimum order and a
minimum-order observer are obtained in $O(\log(\ell+2))$ calls and one lifting step
(Theorem~\ref{thm:interface}), and on a published instance the resulting value is $4$, not the $5$ that
was reported. The feasibility characterization itself is taken from
\cite{fer10,dar25,fer26a,fer26b}.

\subsection{Limitations and what is not claimed}\label{sec:limits}

Table~\ref{tab:models} separates the three models in which the construction is described; the last
row is the one that is open.

\begin{table}[!t]
\centering
\caption{What is claimed, in which model, for the construction of a minimum-order observer.}
\label{tab:models}
\footnotesize
\begin{tabular}{@{}p{0.42\columnwidth}p{0.48\columnwidth}@{}}
\toprule
Task & Claim \\
\midrule
deciding the minimum order & bit complexity, $f(\ell)\operatorname{poly}(\mathcal L)$ \\
finding the algebraic parameters of $W$ & arithmetic complexity, same $f(\ell)$ factor \\
recovering the free rows from $W$ & exact rational lifting, $\operatorname{poly}(\mathcal L)$ \\
the whole construction in the bit model & open: no bit-model sampling-point bound \\
\bottomrule
\end{tabular}
\end{table}

Three qualifications bound what is claimed.

\begin{enumerate}
\item \textbf{The lower bound is not claimed to be tight.} \eqref{eq:qmin-lb} is attained on the
instance of Appendix~\ref{app:example}, where it gives $4$ and $q_{\min}=4$, but the difference
$q_{\min}-\bigl(r_0+\max(0,\operatorname{rank}[U_0;V_0]-s_0+1)\bigr)$ can be positive in general. What
the bound does provide is a test for the orders below it, which then need not be examined.
\item \textbf{The construction is bounded in the arithmetic model.} Corollary~\ref{cor:find} bounds
arithmetic operations and output length, not bit operations: the sampling-point analysis of
\cite{bpr06} that it invokes is arithmetic, and a bit-model sampling-point theorem would be needed
(Remark~\ref{rem:find-bit}).
\item \textbf{The complexity statements are inherited.} Everything quantitative about
$\mathrm{OPT}$ is inherited from the decision procedure of the companion paper through
Theorem~\ref{thm:interface}; this paper does not reprove it, and no closed-form expression for
$q_{\min}$ is offered beyond the characterization of Proposition~\ref{prop:minorder} and the bound of
Proposition~\ref{prop:lowerbound}.
\end{enumerate}

\paragraph*{The size of the parameter}\label{sec:paramsize}

Two quantities describe the free data: the number of coordinates
\begin{equation}\label{eq:kappa-size}
\kappa_{\rm obs}=r\,d_0\le2rn
\end{equation}
needed to write the $r$ free rows, and the parameter that governs the bound,
\begin{equation*}
\ell=\operatorname{rank}(AN)\le k\le n .
\end{equation*}
Only the first can be quadratic in $n$, and only the second bounds the search here: by
Claim~\ref{claim:sandwich} and Claim~\ref{claim:struct} a feasible $W$ satisfies
$0\le\dim(W/W_0)\le2\ell$, so the coordinate count \eqref{eq:kappa-size} records how the free rows are
written, whereas $\ell$ is the quantity that both the minimum-order search and the decision procedure
of the companion paper depend on.

Whether $\ell$ is optimal is open. Two further reductions are conceivable, and they are of different
kinds.

\begin{enumerate}
\item \emph{Quotients of the observer rows.} The conditions depend on row spaces, hence are
invariant under row operations of $L$. Those preserving $L_0$ are block lower-triangular,
\begin{equation*}
T=\begin{bmatrix}I_{r_0}&0\\ U&S\end{bmatrix},\qquad \det S\ne0 ,
\end{equation*}
so two different free blocks can present the same $W$. The reduction of
Section~\ref{sec:standing} already quotients by this symmetry; the question is whether a
smaller invariant than $\ell$ exists.
\item \emph{Spectral invisibility.} One might hope to divide by directions that change neither
the rank condition nor the finite spectrum. Such a quotient must preserve the interaction
between $H$ and $H_A$, not only $\operatorname{rank}H$; a quotient preserving \eqref{eq:5a} but
changing the spectrum would be invalid. We do not know whether one exists.
\end{enumerate}

\subsection{Reproducibility and evidence}\label{sec:repro}

\begin{table*}[!t]
\centering
\caption{Where each statement of the paper is verified. \emph{Literal}: the Lean statement is the
statement of the paper. \emph{Abstract core}: the Lean statement is the linear-map or submodule
version, and the translation to matrices is paper-level. \emph{Ingredient}: one step of the proof,
not the statement. Statements not listed are paper-level throughout.}
\label{tab:evidence}
\footnotesize
\begin{tabular}{@{}p{0.43\textwidth}>{\raggedright\arraybackslash}p{0.15\textwidth}>{\raggedright\arraybackslash}p{0.36\textwidth}@{}}
\toprule
Statement & Lean evidence & Formal name \\
\midrule
Sandwich and realizability (Claim~\ref{claim:sandwich}(i)) & literal &
\texttt{rowSpan\_fromRows\_left}, \texttt{mem\_rowSpan\_iff\_exists\_vecMul} \\
$\pi_1(W)=\pi_2(W)$ from \eqref{eq:5a} and \eqref{eq:5b} at $\lambda=0$
(Claim~\ref{claim:compress}, Claim~\ref{claim:struct}) & abstract core &
\texttt{pi1\_eq\_pi2\_of\_le\_of\_finrank\_eq} \\
Structural bound, upper half (Claim~\ref{claim:struct}), \eqref{eq:struct} & literal &
\texttt{finrank\_le\_two\_mul\_of\_pi1\_eq\_pi2} \\
Structural bound, lower half (Claim~\ref{claim:struct}), \eqref{eq:struct} & literal &
\texttt{finrank\_add\_one\_le\_of\_spectral} \\
Monotonicity (Claim~\ref{lem:monotone}) & abstract core & \texttt{conditions\_fromRows\_zero} \\
Lower bound \eqref{eq:lb}, the inequality $\operatorname{rank}[U_0;V_0]\le\omega$
(Proposition~\ref{prop:lowerbound}; the dimensional half of Proposition~\ref{prop:tight}) &
abstract core & \texttt{finrank\_sup\_pi\_le}, \texttt{finrank\_pi1\_mono} \\
\bottomrule
\end{tabular}
\end{table*}

Two kinds of evidence support the statements above, and the artifact records which statement rests on
which.

\begin{enumerate}
\item \emph{Lean~4 with Mathlib.} The structural statements of Section~\ref{sec:standing}, together
with the dimensional parts of Propositions~\ref{prop:lowerbound} and~\ref{prop:tight}, are formalized;
Table~\ref{tab:evidence} lists them statement by statement, with the grade of the formalization and
the declaration names. The grades separate what is formalized as stated (\emph{literal}) from what is
formalized at the level of linear maps and submodules (\emph{abstract core}), the translation to
matrices being paper-level; the single \emph{ingredient} is one step of the proof of
Lemma~\ref{lem:monotone}. Two points about the coverage. First, the lower half $\omega+1\le\dim W$ of
Claim~\ref{claim:struct} is formalized over an algebraically closed field, so the graph argument that
the spectral condition rules out $\dim W=\dim\pi_1(W)$ for a feasible $W$ is covered. Second,
Proposition~\ref{prop:tight} is only partly covered: what the artifact does cover of it is its
dimensional half, namely, for $W_0\subseteq W$ with $\pi_2(W)\subseteq\pi_1(W)$,
$\dim(\pi_1(W_0)+\pi_2(W_0))\le\dim\pi_1(W)$ and $\dim\pi_1(W_0)\le\dim\pi_1(W)$, which is
$\operatorname{rank}[U_0;V_0]\le\omega$ for $U_0=\pi_1(W_0)$, $V_0=\pi_2(W_0)$. Every declaration is
checked by the kernel against only (\texttt{propext}, \texttt{Classical.choice}, \texttt{Quot.sound}).
The Popov--Belevitch--Hautus criterion that turns $\dim W=\dim\pi_1(W)$ into a kernel condition is
paper-level, since the artifact contains no controllability theory; so are the interface of
Section~\ref{sec:compute} and the counting arguments, hence this is not a formal verification of the
complexity statements.
\item \emph{Exact rational arithmetic.} The instance of Appendix~\ref{app:example} is verified by exact
rational computation, not by numerical evaluation: the two rank equalities of \eqref{eq:5a}, the
vanishing of the determinant of the $8\times8$ pencil together with the nonzero $7\times7$ minors and
their gcd, and the rank of $[U_0;V_0]$ that enters \eqref{eq:qmin-lb}. This excludes arithmetic slips
but does not replace the proofs.
\end{enumerate}

The Lean development and the exact-arithmetic scripts form the artifact at
\url{https://github.com/Zhengshuji/q1-functional-observer}. Of the classical facts invoked by citation,
the rank identity behind Claim~\ref{claim:compress} is also proved in the artifact, where it is used at the
level of linear maps; the remaining cited facts are used as references and are not formalized.

\appendices

\section{A published value corrected}\label{app:example}

We run the procedure on the instance of \cite[Sec.~VII]{fer26b}, which has $n=8$, $p=2$ and
\begin{equation*}
\begin{gathered}
C=\begin{bmatrix}-2&-2&2&-2&-1&0&1&-3\\ 0&0&1&2&2&3&-1&-2\end{bmatrix},\\
L_0=\begin{bmatrix}-2&-1&-2&-3&-1&3&2&0\\ 1&-1&-2&1&0&2&0&1\end{bmatrix},
\end{gathered}
\end{equation*}
and
\begin{equation*}
A=\begin{bmatrix}
0&1&0&1&0&0&0&0\\
0&-1&1&0&0&0&0&-1\\
1&0&-1&1&0&0&1&-1\\
-1&-1&0&-2&0&0&0&0\\
0&0&0&0&-2&0&0&1\\
0&0&0&0&0&-2&1&0\\
0&0&0&0&0&-1&-2&1\\
0&0&0&0&0&0&1&-2
\end{bmatrix}.
\end{equation*}
The compression of Section~\ref{sec:standing} gives
\begin{equation*}
\begin{gathered}
k=\dim\ker[C;CA]=4,\qquad d_0=\operatorname{rank}[N\ AN]=6,\\
s_0=\dim\operatorname{row}(L_0M)=2,
\end{gathered}
\end{equation*}
and, in the notation of \eqref{eq:ell},
\begin{equation*}
S=\operatorname{row}(AN)=\mathbb R^4,\qquad \ell=4,\qquad m'=4 .
\end{equation*}
Here $\ell=k$, so the subspace reduction is vacuous and the content of the example is the
certificate alone. The coordinate count is $\kappa_{\rm obs}=rd_0=12$ already at $r=2$, whereas
the parameter of the companion paper is $\ell=4$.

At $t=2$, that is at order $q=r_0+t=4$, the decision procedure of the companion paper accepts. An explicit
witness is the four-row observer matrix $L=[L_0;R]$ with
\begin{equation*}
R=\begin{bmatrix}
-\tfrac{42919}{32067} & -\tfrac{55096}{10689} & 1 & 0 & 0 & 0 & 0 & 0\\
\tfrac{3283}{4581} & \tfrac{118}{1527} & 0 & 1 & 0 & 0 & 0 & 0
\end{bmatrix}.
\end{equation*}
For this matrix,
\begin{equation*}
\operatorname{rank}[L;CA;C]=7=\operatorname{rank}[LA;L;CA;C],
\end{equation*}
which is \eqref{eq:5a}. The pencil $[\lambda L-LA;\,CA;\,C]$ is $8\times8$; all $64$ of its
$7\times7$ minors are nonzero and have gcd $1$, while its determinant vanishes, so its rank is $7$
for every $\lambda\in\mathbb C$, which is \eqref{eq:5b}. Hence $\mathrm{DEC}(4)$ holds and
$q_{\min}\le4$ for this instance.

We record that \cite[Sec.~VII-B]{fer26b} reports this instance as infeasible at order $4$, and
hence $\{5,6\}$ as its order spectrum. The matrix $L$ above is a four-row witness satisfying the
two conditions, so \emph{that report is in error}. We do not identify the source of the discrepancy:
to do so we would have to reproduce the enumeration used there and locate the branch at fault, which
we have not done. The explicit witness is enough to contradict the reported infeasibility at order
four, and that is all we use. Nothing else in this paper depends on the discrepancy: the
characterization of \cite{fer26b} is used only for the existence of a nullspace matrix, and the
witness supplies one.

\paragraph*{The two procedures on this instance}
Write $M=[N\ AN]$ and, for the accepted $L$, $H=LN$, $H_A=LAN$.

\emph{OPT.} Algorithm~\ref{alg:opt} with $Q=6$ evaluates $\mathrm{DEC}(6)$, which is feasible, so the
answer is finite; the search interval is $[0,\min(4,2\ell)]=[0,4]$. The call at $t=2$, that is at
order $4$, is the certificate above and accepts. The call at $t=0$, that is at $L=L_0$, rejects: for
$L_0$ we have $\operatorname{rank}[L_0;CA;C]=6$ but $\operatorname{rank}[L_0A;L_0;CA;C]=7$, so
\eqref{eq:5a} fails. At $t=1$, that is at order $3$, no feasible $W$ exists, and Proposition~\ref{prop:lowerbound}
decides it without any spectral test: the two blocks of $W_0$ are $U_0=L_0N$ and $V_0=L_0AN$ with
$\operatorname{rank}U_0=2$ and $\operatorname{rank}[\,U_0;\,V_0\,]=3$, so \eqref{eq:qmin-lb} gives
$t\ge3-2+1=2$, i.e.\ every feasible $W$ has quotient dimension at least $2$. (A candidate with $t=1$
would have $\dim W=3$, forcing $\omega=2$ by Claim~\ref{claim:struct} and hence $U=U_0\supseteq V_0$,
which fails.) Therefore $\mathrm{DEC}(3)$ holds for no candidate, $\mathrm{DEC}(4)$ holds by the
certificate above, and

\begin{equation*}
q_{\min}=4
\end{equation*}

for this instance. The value $5$ reported in \cite[Sec.~VII]{fer26b} is therefore not $q_{\min}$.

\emph{FIND.} For the accepted $L$ the subspace $W=\operatorname{row}[H\ H_A]$ is four-dimensional and
$W_0=\operatorname{row}[H_0\ H_{A,0}]$ is two-dimensional, so $t=\dim(W/W_0)=2$; moreover $W_0\subseteq
W$, and the rows $r_1M$ and $r_2M$ of the two free rows, together with the two rows of $W_0$, span
$W$:
\begin{equation*}
\begin{gathered}
\dim W=4,\qquad \dim W_0=2,\\
\operatorname{rank}[\,W_0;\,r_1M;\,r_2M\,]=4 .
\end{gathered}
\end{equation*}
Step~2 of Procedure~\ref{proc:find} is therefore realized by $u_i=r_i$, and step~3 pads with
$Q-r_0-t=6-2-2=2$ zero rows. Only step~1 of the procedure is not exercised here, because a witness
$W$ was found directly rather than by solving the algebraic system.

The two exact rank computations of $U_0$ and $[U_0;V_0]$ are the whole of the argument at order
$3$; they are reproduced by the script archived with the artifact, together with the rank and
minor computations used above.

\begin{thebibliography}{99}
\bibitem{fer26a} J. A. Fernando. A minimum-order functional observer beyond the Darouach and Luenberger constructions. arXiv:2609.33615, 2026 (preprint).
\bibitem{fer26b} J. A. Fernando. Complete characterization of minimum-order functional observers from Darouach to Luenberger. arXiv:2609.37992, 2026 (preprint).
\bibitem{fer10} J. A. Fernando, H. Trinh, L. Jennings. Functional observability and the design of minimum order linear functional observers. \emph{IEEE Trans. Automat. Control} 55(5):1268--1273, 2010.
\bibitem{hj13} R. A. Horn, C. R. Johnson. \emph{Matrix Analysis}. 2nd ed.
Cambridge University Press, 2013.
\bibitem{dar25} M. Darouach, J. A. Fernando. On functional observability and functional observer design. \emph{Automatica}, vol. 173, Art. no. 112115, 2025, doi: 10.1016/j.automatica.2025.112115.
\bibitem{bpr06} S. Basu, R. Pollack, M.-F. Roy. \emph{Algorithms in Real
  Algebraic Geometry}, 2nd ed. Berlin, Germany: Springer, 2006,
  Algorithms and Computation in Mathematics, vol. 10.
\bibitem{lan02} S. Lang. \emph{Algebra}. Revised 3rd ed.
Graduate Texts in Mathematics 211, Springer, 2002.
\bibitem{companion1} Deciding functional observer feasibility: compression, certificates, and
parameterized algorithms. Companion manuscript, 2026. Cited for the decision procedure, its
finite certificate and its complexity bounds.
\end{thebibliography}

\end{document}
